\section{平台、工具与用户心智}

Agent 能做事之后，还要回答它在用户心智里是什么。这个问题决定产品边界：是工具，是平台，还是一个新的入口。

张月光还讨论了新平台机会。他认为新内容平台需要新媒介和新交互；Sora 不一定会成为平台，因为它可能仍然是短视频内容形态；ChatGPT 则可能成为平台，因为它形成了新的交互入口。对 Agent 来说，垂直/通用、模型/应用这些词都太粗，真正要看用户用什么心智记住你。

\begin{figure}[H]
\centering
\includegraphics[width=0.94\textwidth]{platform-tool-agent-positioning.png}
\caption{新平台的机会：工具型产品、平台型产品与 Agent 入口。}
\end{figure}

\begin{knowledgebox}{读图：用户心智决定边界}
工具解决明确任务，平台形成关系和生态，Agent 入口则要让用户为一个高频心智记住品牌。过窄会低频，过宽会“什么都能干等于什么都不能干”。Dokie 的定位是创作与表达 Agent，而不是所有任务的大杂烩。
\end{knowledgebox}

\subsection{创作与表达的边界}

张月光认为，PPT、文档、表达、创作可以合在一个更高层心智里：帮助用户更好地创作和表达。相反，资料收集、填表、信息处理等任务未必会和创作 Agent 合并。这个判断很产品：不是按模型能力分界，而是按用户心智和使用场景分界。

\begin{importantbox}{Agent 的定位问题}
Agent 不是越通用越好。产品要找到一个用户愿意记住、足够高频、足够宽但不空泛的心智边界。
\end{importantbox}

\subsection{本章小结}

Agent 产品的竞争，不只是模型能力竞争，也是用户心智竞争。能否从工具变成入口，取决于它是否拥有高频边界、关系感和持续扩展能力。

